
\documentclass[
12pt,
a4paper
]{article}

% ------------------------------------------------
% Pakete
% ------------------------------------------------

\usepackage[utf8]{inputenc}
\usepackage[T1]{fontenc}
\usepackage[german]{babel}

\usepackage[a4paper,margin=2.5cm]{geometry}
\usepackage{graphicx}
\usepackage{booktabs}
\usepackage{longtable}
\usepackage{float}
\usepackage{parskip}
\usepackage[hidelinks]{hyperref}

% ------------------------------------------------
% Dokumentinformationen
% ------------------------------------------------

\title{
\Huge Ergebnisdokumentation\\[0.5cm]
\Large Workflowbasierte Rückbestätigung von Lieferterminen
}

\author{
THWS \\[0.2cm]
in Kooperation mit MINOVA
}

\date{\today}

\begin{document}

% =================================================
% TITELSEITE
% =================================================

\begin{titlepage}
\centering

\vspace*{3cm}

{\Huge\bfseries Ergebnisdokumentation\par}

\vspace{1cm}

{\LARGE Workflowbasierte Rückbestätigung von Lieferterminen\par}

\vspace{2cm}

{\Large Projektarbeit\par}

\vspace{1cm}

\begin{tabular}{ll}
Unternehmen: & MINOVA GmbH \\
Hochschule: & THWS \\
Projektzeitraum: & xx.xx.2026 -- xx.xx.2026 \\
Team: & Name 1, Name 2, Name 3 \\
\end{tabular}

\vfill

Stand: \today

\end{titlepage}

\tableofcontents
\newpage

% =================================================
\section{Einleitung}
% =================================================

\subsection{Projektbeschreibung}

Im Rahmen der Projektarbeit wurde in Kooperation mit der
MINOVA GmbH ein System zur workflowbasierten Rückbestätigung
von Lieferterminen entwickelt.

Ziel des Projekts war die Digitalisierung und Automatisierung
der Kommunikation zwischen Unternehmen und Kunden bei der
Abstimmung von Lieferterminen.

\subsection{Ziel der Ergebnisdokumentation}

Diese Ergebnisdokumentation beschreibt die Ausgangssituation,
die angestrebten Ergebnisse, die umgesetzte Lösung sowie die
fachlichen und technischen Erkenntnisse des Projekts. Der Fokus
liegt dabei auf dem erreichten Nutzen und nicht auf der
Projektorganisation oder der zeitlichen Planung.

% =================================================
\section{Ausgangssituation}
% =================================================

\subsection{IST-Situation}

Vor Projektbeginn erfolgte die Abstimmung von Lieferterminen
größtenteils manuell. Kunden wurden über geplante Lieferungen
informiert und mussten Rückmeldungen auf unterschiedlichen
Kommunikationswegen übermitteln.

Die Informationen wurden anschließend manuell verarbeitet und
an die beteiligten Stellen weitergegeben.

\subsection{Fachlicher Kontext}

Die Rückbestätigung von Lieferterminen ist ein zentraler
Bestandteil der operativen Abstimmung zwischen Unternehmen,
Disposition und Kunden. Verzögerungen oder unvollständige
Rückmeldungen führen unmittelbar zu Mehraufwand in der
Kommunikation und erschweren eine verlässliche Planung.

Vor diesem Hintergrund bestand die Anforderung darin, einen
nachvollziehbaren und digital unterstützten Prozess zu schaffen,
der Rückmeldungen strukturiert erfasst und ohne zusätzliche
manuelle Zwischenschritte weiterverarbeitet.

\subsection{Identifizierte Probleme}

\begin{itemize}
    \item Hoher Kommunikationsaufwand
    \item Manuelle Bearbeitung von Rückmeldungen
    \item Medienbrüche zwischen verschiedenen Systemen
    \item Eingeschränkte Transparenz über den aktuellen Status
    \item Erhöhter organisatorischer Aufwand
\end{itemize}

% =================================================
\section{Ergebnisziele}
% =================================================

\subsection{Fachliche Ziele}

\begin{itemize}
    \item Automatisierte Kundenkommunikation
    \item Digitale Rückbestätigung von Lieferterminen
    \item Transparente Statusverwaltung
    \item Verbesserte Nachvollziehbarkeit
\end{itemize}

\subsection{Technische Ziele}

\begin{itemize}
    \item Entwicklung einer Webanwendung
    \item Integration von Frontend und Backend
    \item Workflowsteuerung über Camunda
    \item Automatisierter E-Mail-Versand
    \item Speicherung relevanter Informationen in einer Datenbank
\end{itemize}

\subsection{Abgrenzung}

Die Ergebnisdokumentation betrachtet insbesondere die
funktionsbezogenen Resultate des Projekts. Themen wie
Projektaufwand, Zeitplanung, Arbeitsverteilung und die
teambezogene Reflexion werden in der separaten
Projektdokumentation behandelt.

% =================================================
\section{Umsetzung}
% =================================================

\subsection{Systemübersicht}

Die entwickelte Lösung besteht aus mehreren Komponenten:

\begin{itemize}
    \item React-Frontend
    \item Spring-Boot-Backend
    \item Camunda Workflow Engine
    \item PostgreSQL-Datenbank
    \item E-Mail-Service
\end{itemize}

\subsection{Prozessablauf}

Der umgesetzte Prozess läuft wie folgt ab:

\begin{enumerate}
    \item Ein Lieferauftrag wird verarbeitet.
    \item Der Kunde erhält automatisch eine E-Mail.
    \item Der Kunde bestätigt oder lehnt das Zeitfenster ab.
    \item Die Rückmeldung wird automatisch verarbeitet.
    \item Der Status wird aktualisiert.
    \item Die Informationen werden im System dargestellt.
\end{enumerate}

% =================================================
\section{Erzielte Ergebnisse}
% =================================================

\subsection{Erreichte Funktionen}

Im Rahmen der Projektarbeit konnten folgende Ergebnisse erzielt
werden:

\begin{itemize}
    \item Funktionsfähige Kommunikation zwischen Frontend und Backend
    \item Automatisierter Versand von E-Mails
    \item Verarbeitung von Bestätigungen und Ablehnungen
    \item Automatische Weiterleitung von Informationen
    \item Darstellung der Informationen auf den Übersichtsseiten
    \item Integration der Workflowsteuerung
\end{itemize}

\subsection{Validierung}

Das entwickelte System wurde gemeinsam mit dem Auftraggeber
vorgestellt und getestet.

Dabei konnte bestätigt werden, dass:

\begin{itemize}
    \item die Kommunikation zwischen Frontend und Backend funktioniert,
    \item Rückmeldungen korrekt verarbeitet werden,
    \item Statusänderungen ordnungsgemäß dargestellt werden,
    \item die Benutzerführung nachvollziehbar ist.
\end{itemize}

\subsection{Abgleich mit den Ergebniszielen}

Die zu Beginn definierten fachlichen und technischen Ziele
konnten in wesentlichen Punkten erreicht werden. Insbesondere
wurden die digitale Rückbestätigung von Lieferterminen, die
strukturierte Statusverwaltung sowie die workflowgestützte
Verarbeitung erfolgreich umgesetzt.

Darüber hinaus konnte gezeigt werden, dass die gewählte
Architektur eine tragfähige Grundlage für weitere
Ausbaustufen bietet. Die erzielten Ergebnisse gehen damit
über einen rein konzeptionellen Entwurf hinaus und bilden
einen praktisch nutzbaren Prototypen ab.

% =================================================
\section{Ausblick}
% =================================================

Für zukünftige Erweiterungen wurden folgende Ansätze
identifiziert:

\begin{itemize}
    \item Kartenansicht zur Anzeige der Fahrerposition
    \item Erweiterung der E-Mail-Kommunikation
    \item Zusätzliche Automatisierungen
    \item Integration weiterer Systeme
\end{itemize}

% =================================================
\section{Erfahrungen und Erkenntnisse}
% =================================================

Im Verlauf der Umsetzung zeigte sich, dass insbesondere die
klare Abbildung fachlicher Entscheidungen in einem definierten
Workflow einen hohen Mehrwert für die Nachvollziehbarkeit des
Gesamtprozesses bietet. Wiederkehrende Kommunikationsschritte
lassen sich dadurch konsistenter und transparenter gestalten.

Gleichzeitig wurde deutlich, dass die Abstimmung zwischen
Benutzeroberfläche, Backend-Logik und Workflowsteuerung
frühzeitig präzisiert werden muss. Schnittstellen,
Statusübergänge und E-Mail-Inhalte wirken direkt auf die
Benutzererfahrung und sollten deshalb fachlich und technisch
gemeinsam betrachtet werden.

Als wesentliche Erkenntnis lässt sich festhalten, dass schon
eine gezielte Teilautomatisierung spürbare Verbesserungen in
Bezug auf Transparenz, Bearbeitungsgeschwindigkeit und
Nachverfolgbarkeit erzeugen kann.

% =================================================
\section{Fazit}
% =================================================

Die definierten Projektziele konnten erfolgreich umgesetzt
werden. Die entwickelte Lösung ermöglicht eine automatisierte
und transparente Rückbestätigung von Lieferterminen und trägt
zur Reduzierung manueller Arbeitsschritte bei.

Die erfolgreiche Integration von Frontend, Backend und
Workflow Engine bildet eine solide Grundlage für zukünftige
Erweiterungen und den weiteren Ausbau des Systems.

\end{document}
